\documentclass[11pt,a4paper]{article}
\usepackage{../shared/luxcover}   % black cover (self-contained, TikZ logo, no assets)
\usepackage{../shared/proposal}   % \proposalhead helpers, stake environment, muted colour
\usepackage[utf8]{inputenc}\usepackage[T1]{fontenc}
\usepackage{amsmath,amssymb}
\usepackage{listings}\usepackage{xcolor}\input{../shared/lstlang}

\newcommand{\code}[1]{\texttt{\small #1}}
\newcommand{\prov}[1]{{\color{muted}\scriptsize\sffamily[#1]}}
\setlength{\emergencystretch}{3em}

\title{Agents That Cost Bytes}
\author{Zach Kelling\\\textit{Lux Industries, Inc.}\\\texttt{research@lux.network}}
\date{2026}

\begin{document}
\luxcoverpage
\sloppy

\noindent{\color{muted}\footnotesize\sffamily ADOPT-14 \quad\textbullet\quad Hanzo AI \& Lux Industries \quad\textbullet\quad 2026}\par
\vspace{0.9ex}

\noindent\textbf{The number.}
An agent costs 601 bytes of heap while it is live and 477 bytes on disk while it is not
\prov{measured}. A million live agents occupy 573.0\,MB, spawn in 125\,ns each, and wake in
107\,ms --- 107\,ns apiece \prov{measured}. A million dormant agents occupy 455\,MB and are
written at 309,789 per second; resuming one in process costs 0.034\,ms \prov{measured}. These
are one pass of \code{hanzoai/cloud bench} on one laptop, and the host is named beside them
because an earlier version of this table was not reproducible: re-running the goroutine lane on
a slower M1 Max returned 452\,ns to spawn against a published 187 and 347\,ns to wake against
114 \prov{verified}. The figures stand; the machine is now part of the figure.

\vspace{0.35ex}
\noindent\textbf{The problem, in their numbers.}
Morpheus sells agents and does not run any. Seven repositories in the organisation carry
``agent'' in the name \prov{verified 2026-09-17}. Four --- \code{morpheus-open-agents},
\code{morpheus-coding-agent}, \code{morpheus-community-agent},
\code{morpheus-knowledge-agent} --- are forks of \code{vercel-labs} templates, every one of
them last touched 2026-04-21 \prov{verified}. Two more are dormant:
\code{agent-computer-use} (2025-07-22) and \code{SmartAgents}, itself a fork, last touched
2024-05-28 \prov{verified}. The seventh, \code{MySuperAgent}, is the live one (2026-09-08 in
our clone) and is a Next.js web application: its own documentation describes ``client-side
agent orchestration and execution'' over ``local SQLite for client-side data storage''
(\code{CLAUDE.md:17--18}) \prov{verified}. None of the seven is a server-side agent runtime.

The sandbox is the sharper case. \code{MySuperAgent} exposes a tool whose description reads
\emph{``Execute Python code in a safe sandbox environment''} (\code{code-tools.ts:6}). The
implementation executes nothing. It returns a string, above a comment reading ``In a real
implementation, you'd use a sandbox like Pyodide or a secure container''
(\code{code-tools.ts:17--19}) \prov{verified}. An agent marketplace whose code executor is a
stub has not deferred the hard part; it has not reached it. Nothing settles on work that never
ran, which is the same meter problem ADOPT-07 names, arriving one layer up.

\vspace{0.35ex}
\noindent\textbf{What a comparison is worth.}
Numbers against competitors are reported here in three columns, because two is how a wrong one
survives. \emph{Ours} is what this pass measured. \emph{Theirs} is a figure attributed to
somebody else and not re-run --- their protocol, their hardware, their claim. \emph{Checked} is
what their library did when we ran it here.

The third column is the one that matters. This lane previously reported a V8 context at
0.15\,ms against a cited 2.79\,ms and called the difference seventeenfold. Those are not the
same primitive: \code{vm.createContext} makes a fresh global inside an isolate that is already
running, and it is cheap precisely because it shares the heap --- the one thing an isolate does
not do. Running the library the 2.79\,ms is attributed to, \code{isolated-vm} 7.0.1, on this
laptop, an isolate costs 0.35\,ms and 1.00\,MiB \prov{checked}. Still under the figure cited at
us, by 4.7$\times$ rather than 17$\times$, and a measurement instead of an argument. The
upstream commit that corrected it is titled ``Measure the isolate instead of citing one''.

So the megabyte is V8's, not a vendor's: any isolate costs about that, including one of ours.
The honest comparison is 601 bytes for a live agent against 1.00\,MiB for an isolate
\prov{checked}, and a 125\,ns spawn against 0.35\,ms \prov{checked}. A comparison is worth
printing only when both sides are the same primitive.

The same discipline cuts against us elsewhere, and is printed anyway. \code{hanzo-vm} cold
boots in 309\,ms and resumes from a checkpoint in 311\,ms \prov{measured}, against E2B under
200\,ms, Modal near a second, and Cloudflare one to three seconds \prov{attributed}. On cold
boot E2B is ahead. The density claim is about agents, not about virtual machines, and stating
where it does not hold is what makes the rest checkable.

\vspace{0.35ex}
\noindent\textbf{The fix, named to a repo.}
A Morpheus provider runs \code{Morpheus-Lumerin-Node/proxy-router} and nothing in it changes.
Alongside it, \code{hanzoai/cloud} supplies the runtime the agent repositories assume and do
not contain: agents that persist as bytes rather than processes, an execution surface that
actually executes, and a meter that counts what ran. Extensions host in process from the same
pass --- \code{wazero} instantiates a WASM module in 7.8\,\textmu s and calls into it in
20\,ns; \code{goja} costs 2.4\,\textmu s per JavaScript VM and 740\,ns per warm eval;
\code{gpython} costs 25.0\,\textmu s per context \prov{measured}. An agent that needs a
language runtime does not need a container for it.

The density is the economic argument. At 601 bytes live and 477 dormant, a provider hosts
agent populations per box that a per-isolate or per-container model prices out of reach: the
isolate at 1.00\,MiB \prov{checked} is roughly 1,700$\times$ the live cost, before any
container overhead. Capacity, work and identity meter as three quantities, and work pays only
on counts the consumer also signed --- so a faked session earns zero, and the settlement
problem ADOPT-07 describes for inference is solved the same way for agents.

\vspace{0.35ex}
\noindent\textbf{The outcome, projected.}
Assumptions stated. Morpheus keeps its contracts, its node, and its token. It gains a runtime
under the seven agent repositories that currently have none, a code executor that is not a
stub, and a meter that makes agent work settleable on the same rails as inference. The
reciprocal case is ADOPT-10, and the open-source path ADOPT-15: the runtime is open, so a provider who adopts it reads every
line of what meters them. What Morpheus supplies in return is demand and a settlement layer
already live on Base --- which is the loop, not a favour.

\vspace{0.35ex}
\noindent\textbf{What would falsify this.}
A published Morpheus agent runtime that executes untrusted code server-side, with its own
density and isolation figures on a named host. None was found across the seven agent
repositories on 2026-09-17, when the organisation stood at 63 public repositories
\prov{verified}. That is a search of the named surface, not a proof of absence across all 63.
If one ships, this brief should be re-run against it rather than restated.

\end{document}
